% GENERATED by build_do_now_tex.py on 2026-04-24 --- do not hand-edit; regenerate instead.
\documentclass[11pt]{article}
\usepackage{preamble}

\begin{document}

\packetheading{Lesson 5-5 \textperiodcentered\ Period 1 \textperiodcentered\ Do Now}{L55 P1}

\vspace{0.25em}
\noindent\textbf{Name:}\ \rule{2in}{0.4pt}\quad
\textbf{Period:}\ \rule{0.5in}{0.4pt}\quad
\textbf{Date:}\ \rule{1in}{0.4pt}
\vspace{0.5em}

\bankitem{Do Now}{%
In business, the term profit is used to describe the difference between the money the business earns (revenue) and the money the business spends (cost).

[IMAGE: Illustration of a man grooming a bulldog, with a sign reading "GROOMING $25 per pet".]

A. Let x represent the number of pets groomed in a month at Grooming USA. Define a revenue function for the business.

B. Materials and labor for each pet groomed cost $15. The business also has fixed costs of $1,000 each month. Define a cost function for this business.

C. Last month, Grooming USA groomed 95 pets. Did they earn a profit? What would the profit be if the business groomed 110 pets in a month?

D. Generalize Explain your procedure for calculating the profit for Grooming USA. Suppose you wanted to calculate the profit for several different scenarios. How could you simplify your process?
}{}
\vspace{2.5in}

\end{document}
